\chapter{Elastizität und sequenzielle Thermomechanik}\label{chap:mechanics}
\section{Kleine Verzerrungen und virtuelle Arbeit}
Unter quasistatischen Bedingungen ist
\begin{equation}\label{eq:mechanical-strong}
-\nabla\cdot\boldsymbol\sigma=\mathbf b,\qquad
\boldsymbol\sigma\mathbf n=\bar{\mathbf t}\text{ auf }\Gamma_t,
\qquad\mathbf u=\bar{\mathbf u}\text{ auf }\Gamma_u.
\end{equation}
Für $\mathbf v=0$ auf $\Gamma_u$ folgt mit der Symmetrie von $\sigma$
\begin{equation}\label{eq:mechanical-weak}
\int_\Omega\varepsilon_v(\mathbf v)^T\sigma_v\,\mathrm dV
=\int_\Omega\mathbf v^T\mathbf b\,\mathrm dV+
\int_{\Gamma_t}\mathbf v^T\bar{\mathbf t}\,\mathrm dA.
\end{equation}
Die Identität $\nabla\mathbf v:\sigma=\operatorname{sym}\nabla\mathbf v:\sigma$
entfernt den antisymmetrischen Anteil. Ein isotropes Material erfüllt
$\sigma=2\mu\varepsilon+\lambda\operatorname{tr}(\varepsilon)\mathbf I$ mit
\begin{equation}\label{eq:lame}
\lambda=\frac{E_Y\nu}{(1+\nu)(1-2\nu)},\qquad
\mu=\frac{E_Y}{2(1+\nu)}.
\end{equation}
In der Engineering-Konvention sind die Normaldiagonalen $\lambda+2\mu$,
die Normalkopplungen $\lambda$ und die drei Schubdiagonalen $\mu$, nicht $2\mu$.
Anisotrope Daten werden als globale SPD-Matrix $\mathbf C$ eingegeben.
Isotrope Eingaben benötigen $E_Y>0$, $-1<\nu<0{,}49$; das strengere obere Limit
als die mathematische Grenze $0{,}5$ schützt die freigegebene
Verschiebungsformulierung. Tet10 ist keine automatische Lösung des
Inkompressibilitätslockings. Körperkräfte tragen N/m$^3$, Traktionen Pa.
\index{Elastizität}\index{Starrkörpermode}

\section{Diskrete Steifigkeit, Lasten und Lagerung}
Mit der Matrix aus Gleichung~\ref{eq:B} folgen
\begin{equation}\label{eq:mechanical-element}
\mathbf K_u^e=\int_K\mathbf B^T\mathbf C\mathbf B\,\mathrm dV,
\quad\mathbf f_b^e=\int_K(\mathbf N\otimes\mathbf I_3)^T\mathbf b\,\mathrm dV,
\quad\mathbf f_t^F=\int_F(\mathbf N\otimes\mathbf I_3)^T\bar{\mathbf t}\,\mathrm dA.
\end{equation}
Tet4 liefert exakt $V\mathbf B^T\mathbf C\mathbf B$ und
$V\mathbf b/4$ je Knoten für konstante Lasten. Globale Steifigkeit hat N/m,
Lasten N. Punktlasten sind integrierte Kräfte auf ausgewählten DOFs; sie sind
keine räumlich aufgelöste Kontakt- oder Spannungssingularitätsmodellierung.

Ein unbelasteter freier Körper besitzt drei Translationen und drei Rotationen
$\mathbf u=\mathbf c+\boldsymbol\omega\times\mathbf x$ mit $\varepsilon=0$.
Die Lagerung muss alle verbleibenden Starrkörpermoden je Komponente entfernen;
ein bloßes Zählen von sechs festgesetzten DOFs genügt nicht.
Die Rangprüfung bezieht sich auf die geometrisch aufgebauten Moden.
Ein Tet4 hat Steifigkeitsrang sechs bei zwölf DOFs, ein Tet10 Rang 24 bei 30
DOFs für ein auflösbares positives Material. Globale Mechanismen und schlechte
Kondition können trotzdem zusätzliche Diagnostik verlangen.

Gespeicherte elastische Energie ist $U_m=\frac12\int\varepsilon_v^T\mathbf C
\varepsilon_v$. Reaktionen sind die ursprünglichen Residuen auf festen DOFs.
Für statisches Gleichgewicht schließen die Summe äußerer Kräfte und Reaktionen
sowie deren Momente $\sum\mathbf x_i\times\mathbf F_i$.
Ein kleines algebraisches Residuum ersetzt diese zwei physikalischen Bilanzen
nicht. Die Reaktionsauswertung erfolgt vor einer Interpretation als
Spannungsresultierende, nicht auf dem mit Einheitszeilen modifizierten System.

\subsection{Handrechnung: affiner geklemmter Dehnungszustand}
Im Einheits-Tet4 sei $E_Y=100$ Pa, $\nu=0{,}25$, also $\lambda=\mu=40$ Pa.
$\mathbf u=(0{,}01x,0,0)$ m ergibt
$\varepsilon_v=(0{,}01,0,0,0,0,0)$ und
$\sigma_v=(1{,}2,0{,}4,0{,}4,0,0,0)$ Pa.
Die transversalen Spannungen entstehen wegen der vorgegebenen Null-
Querdehnungen; dies ist kein frei kontrahierender Zugstab.
$U_m=(1/6)(0{,}01\cdot1{,}2)/2=10^{-3}$ J.
Bei einem \emph{frei} kontrahierenden Zugstab wäre dagegen
$u_x=Fx/(AE_Y)$ und $\varepsilon_{yy}=\varepsilon_{zz}=-\nu\varepsilon_{xx}$.
Die Randbedingungen entscheiden, welche Referenz anwendbar ist.

\section{Thermische Eigenverzerrung und Kopplungsrichtung}
Für bekannte Temperatur lautet
\begin{equation}\label{eq:thermal-strain}
\varepsilon_{\rm th}=\alpha(T-T_{\rm ref}),\qquad
\varepsilon_m=\varepsilon_v-\varepsilon_{\rm th},\qquad
\sigma_v=\mathbf C\varepsilon_m.
\end{equation}
In Gleichung~\ref{eq:mechanical-weak} eingesetzt und den bekannten Anteil
nach rechts verschoben ergibt
\begin{equation}\label{eq:thermal-load}
\mathbf K_u\mathbf u=\mathbf f_b+\mathbf f_t+\mathbf f_{\rm th},\qquad
\mathbf f_{\rm th}^e=\int_K\mathbf B^T\mathbf C\alpha
(\mathbf N\mathbf T_e-T_{{\rm ref},K})\,\mathrm dV.
\end{equation}
$\alpha$ und $T_{{\rm ref},K}$ dürfen regionsweise verschieden sein. Die
Temperatur wird auf demselben Netz zuerst gelöst; Mechanik speist keine
Wärmeleistung zurück. Es gibt keine reversible Entropie-/Trägheits-
Thermokopplung und keine Wärmeproduktion durch mechanische Dissipation.
\index{Thermomechanik}\index{Eigenverzerrung}

Bei Tet4 sind $\mathbf B,\mathbf C,\alpha$ konstant. Daher ist die Last exakt
mit $\bar T_K=\sum_a T_a/4$ und
$\mathbf f_{\rm th}^e=V\mathbf B^T\mathbf C\alpha(\bar T_K-T_{{\rm ref},K})$.
Der berichtete Tet4-Spannungszustand und die mechanische Zellenergie verwenden
auch diese mittlere Eigenverzerrung. Für räumlich veränderliches $T$ ist die
Lastmittelung exakt, die berichtete Energie aber die Energie dieses diskreten
Zellmittelmodells. Die kontinuierliche Energie mit punktweiser Eigenverzerrung
enthielte zusätzlich den temperaturabhängigen Varianzterm
\begin{equation}\label{eq:thermal-variance}
U_{m,\rm cont}-U_{m,\rm mean}
=\frac V2\left(\overline{\varepsilon_{\rm th}^T\mathbf C\varepsilon_{\rm th}}
-\bar\varepsilon_{\rm th}^T\mathbf C\bar\varepsilon_{\rm th}\right)\ge0.
\end{equation}
Dieser bekannte, von $\mathbf u$ unabhängige Term verändert die mechanische
Gleichung nicht. Er wird nicht still als Bestandteil der Tet4-Ausgabe behauptet.
Tet10 integriert Last und Energie mit $T_h$ an jedem Quadraturpunkt; die
berichtete Spannung ist eine gesonderte Schwerpunktprobe.

\subsection{Handrechnung: frei und vollständig behindert}
Mit denselben elastischen Zahlen, $\alpha=(0{,}002,0,0,0,0,0)$/K und
$T-T_{\rm ref}=10$ K gilt $\varepsilon_{\rm th,x}=0{,}02$.
Bei frei möglicher Expansion ergibt $u_x=0{,}02x$ m, $\sigma=0$, $U_m=0$.
Für vollständig vorgegebene $\mathbf u=0$ sind
$\sigma_v=(-2{,}4,-0{,}8,-0{,}8,0,0,0)$ Pa und
$U_m=(1/6)(0{,}02\cdot2{,}4)/2=0{,}004$ J.
Der thermische Lastvektor ist positiv auf der rechten Seite, obwohl die
behinderte Spannung negativ ist. Diese Vorzeichenprüfung verhindert die
häufige Verwechslung zwischen Eigenverzerrung und äußerer Zuglast.

\subsection*{Implementierung und unabhängige Nachweise.}
{\small\raggedright
\path{src/fem_toolbox/solvers/elasticity.py},
\path{src/fem_toolbox/solvers/thermomechanics.py},
\path{src/fem_toolbox/elements/quadratic.py};
\path{tests/patch/test_elasticity_patch.py},
\path{tests/patch/test_thermomechanics_patch.py},
\path{tests/unit/test_thermal_expansion.py}.
Die gedrehte thermische Expansion wird in einem unabhängigen Patch geprüft;
Doc/14 beschreibt die gekoppelte Manufactured Solution auf vier Netzen.
Tet10-Nachweise: \path{tests/convergence/test_tet10_elasticity.py} und
\path{tests/convergence/test_tet10_thermomechanics.py}.

\par}
